\chapter{回到具体生物学：乳糖操纵子中的有效反馈}
至此，最小模型已经解释了记忆、振荡、适应与随机涨落各自可能来自什么机制。把它们用于真实系统时，第一步是识别生物过程怎样对应到模型变量和调控函数。乳糖操纵子提供了一个例子：反馈并不只发生在基因之间，还经过摄取与代谢。
大肠杆菌 lac 系统包括 LacI 阻遏蛋白、lacZYA 转录单位以及乳糖/异乳糖相关代谢与摄取过程。lacI 具有自身的转录调控，不属于 lacZYA 同一多顺反子转录单位。异乳糖降低 LacI 的阻遏活性；CAP--cAMP 及葡萄糖相关效应调节转录和摄取。

在最简单的结合占据近似下，活性 LacI 浓度为 $R_a$ 时，解除阻遏的分数可写作 $K_d/(K_d+R_a)$，从而给出有效转录率 $F(R_a)$。进一步引入 LacY 蛋白 $y$ 与细胞内诱导物 $a$，
\begin{equation}\dot y=F(a)-\gamma_y y,\qquad\dot a=v_{\rm in}(y,L_{\rm ext})-v_{\rm out/met}(a,y).\end{equation}
这体现了 ``LacY 增强摄取 $\to$ 诱导物增加 $\to$ 阻遏解除 $\to$ LacY 增加'' 的有效正反馈，但真正的乳糖代谢、透过酶动力学、葡萄糖作用和生长稀释会改变反馈强度。特别应区分代谢性乳糖与非代谢型诱导物的实验情形，不能笼统声称 lac 网络必然双稳态。

\begin{figure}[H]\centering
\begin{tikzpicture}
\node[gene] (y) at (0,0) {LacY};
\node[gene] (a) at (3,0) {胞内诱导物};
\node[gene] (r) at (6,0) {活性 LacI};
\node[gene] (f) at (9,0) {lacZYA 表达};
\draw[act] (y)--node[above]{摄取}(a);
\draw[rep] (a)--(r);\draw[rep] (r)--(f);
\draw[act] (f.south) -- ++(0,-.7) -| (y.south);
\end{tikzpicture}
\caption{lac 系统的简化有效反馈关系。两次抑制形成解除阻遏，闭环总体为正；该图省略了代谢、葡萄糖作用和稀释，不据此断言双稳态。}
\end{figure}

为了把这条生物学路径与前面的符号分析接上，可以逐步计算导数。单个位点近似下 $F(R_a)=r_{\rm leak}+(r_{\rm max}-r_{\rm leak})K_d/(K_d+R_a)$，当 $r_{\rm max}>r_{\rm leak}$ 时 $\partial F/\partial R_a<0$。
若有效活性阻遏蛋白 $R_a(a)$ 随诱导物下降，则链式法则给出 $\dd F/\dd a=(\partial F/\partial R_a)(\dd R_a/\dd a)>0$。
再结合 $\partial v_{\rm in}/\partial y>0$，即可识别摄取支路的正反馈。完整第二个方程的交叉导数却是
$\partial\dot a/\partial y=\partial_yv_{\rm in}-\partial_yv_{\rm out/met}$，其符号未必正；这正是代谢性诱导物与仅摄取的理想模型之间需要检查的环节。

这个例子把最小回路的用法具体化了：先识别产生率与中间变量，再检查有效反馈的符号和强度，最后分析对应参数下的动力学。接下来还要问，多个因子共同作用于启动子时怎样合并？局部回路受到其余网络影响时，独立分析还能保留多少结论？
